\problemname{Talia's Supplies}


Talia is preparing to embark on a brave journey to the long-lost island of Numbertheorious,
where the legendary proof of the O(1) graph algorithm is said to be hidden. But before setting
off on her journey, Talia must stop at an expedition warehouse to collect the necessary supplies.
Along one long shelf are N supply crates, arranged from left to right. Each crate has an integer
value that represents its usefulness, depending on how much it will help or hinder her journey.

A positive usefulness value represents supplies that will help Talia, while a negative usefulness
value represents supplies that may make her journey more difficult.

There is one problem. The warehouse uses an old collection system that can only release one
continuous section of the shelf at a time. This means that Talia has to choose exactly one
non-empty section and take every crate between her chosen starting and ending positions.

Talia wants to leave with the greatest possible total usefulness.

Determine the maximum total usefulness Talia can obtain, as well as the positions of the first
and last crates she should collect.

If several sections have the same maximum total usefulness, Talia chooses the section with the
earliest starting position. If there is still a tie, she chooses the section with the earliest
ending position.

Input

The first line contains a single integer N, representing the number of supply crates.
N will be between 1 and 200,000.

The second line contains N integers representing the usefulness values of the crates from left
to right. Each usefulness value will be between -10,000 and 10,000.

Output

Output three integers separated by spaces: the maximum total usefulness, the position of the
first crate selected, and the position of the last crate selected.

Crate positions are numbered from 1 to N.

Explanation

For example, consider the usefulness values:

-4, 6, -2, 7, -8, 3, 5, -6

The best section is crates 2 through 4:

6, -2, 7

The total usefulness is:

6 + (-2) + 7 = 11

Therefore, the output is:

11 2 4